% !TeX program = pdflatex
%==============================================================================
%  Lista Teórica 10 --- Máquinas de Vetores de Suporte (SVM)
%  O gabarito sai de "Lista teorica 09 - gabarito.tex".
%==============================================================================
\providecommand{\gabaritoopt}{}
\documentclass[11pt,a4paper]{article}
\usepackage[\gabaritoopt]{../../recursos/latex/estilo-lista}

\begin{document}
\cabecalholista{09}{Máquinas de Vetores de Suporte (SVM)}

%------------------------------------------------------------------------------
\begin{exercicio}[a margem, na mão]
Quatro observações em $\R^2$, linearmente separáveis:
\begin{center}\small
\renewcommand{\arraystretch}{1.2}
\begin{tabular}{@{}lcccc@{}}
\toprule
        & $A$ & $B$ & $C$ & $D$ \\
\midrule
$\x_i$  & $(3,1)$ & $(3,3)$ & $(1,1)$ & $(1,3)$ \\
$y_i$   & $+1$ & $+1$ & $-1$ & $-1$ \\
\bottomrule
\end{tabular}
\end{center}
Lembre que o hiperplano de margem máxima é escrito na forma canônica
$\bm w^\top\x + b = 0$ com $y_i(\bm w^\top\x_i+b)\ge 1$ para todo $i$, e que a
margem de cada lado vale $1/\norm{\bm w}$.
\begin{enumerate}[label=(\alph*),leftmargin=1.1cm]
  \item Desenhe os quatro pontos e determine, por inspeção, o hiperplano de
        margem máxima. Escreva $\bm w$ e $b$ na forma canônica e calcule a margem.
  \item Quais pontos são vetores de suporte?
  \item Acrescente $E=(5,2)$ com $y_E=+1$. O hiperplano muda? Justifique sem
        refazer a conta.
  \item Agora mova $A$ de $(3,1)$ para $(2{,}5\,;\,1)$, mantendo os demais. Qual
        é o novo hiperplano e a nova margem? Quais são agora os vetores de
        suporte?
\end{enumerate}
\end{exercicio}

\begin{solucao}
\textbf{(a)} Os positivos estão todos em $x_1=3$ e os negativos em $x_1=1$; a
coordenada $x_2$ não distingue nada. O hiperplano que mais se afasta dos dois
grupos é a reta vertical no meio, $x_1=2$.

Na forma canônica precisamos de $y_i(\bm w^\top\x_i+b)=1$ nos pontos mais
próximos. Tomando $\bm w=(1,0)$ e $b=-2$:
\[
  A:\ (1)(3-2)=1,\quad B:\ (1)(3-2)=1,\quad
  C:\ (-1)(1-2)=1,\quad D:\ (-1)(1-2)=1 .
\]
Todos valem exatamente $1$, como exige a forma canônica. A margem é
\[
  \frac{1}{\norm{\bm w}} = \frac{1}{1} = \mathbf{1}
\]
de cada lado, ou seja uma faixa de largura total $2$, entre $x_1=1$ e $x_1=3$.

\smallskip
\textbf{(b)} \textbf{Todos os quatro.} Vetor de suporte é o ponto que satisfaz
$y_i(\bm w^\top\x_i+b)=1$, isto é, que \emph{toca} a margem --- e aqui os quatro
tocam.

\begin{atencao}[o que o \texttt{scikit-learn} vai dizer]
Se você conferir com \texttt{SVC(kernel="linear", C=1e6)} e imprimir
\texttt{support\_vectors\_}, vão aparecer \textbf{dois} pontos, não quatro. Não é
erro: nesta configuração simétrica a solução \emph{dual} não é única, e o
otimizador devolve um conjunto mínimo de multiplicadores que já reproduz o mesmo
$\bm w$. A definição geométrica (pontos sobre a margem) e a computacional
(pontos com multiplicador de Lagrange não nulo) coincidem no caso genérico e
podem divergir em casos degenerados como este.
\end{atencao}

\smallskip
\textbf{(c)} \textbf{Não muda.} Para $E=(5,2)$ temos
$y_E(\bm w^\top\x_E+b) = 5-2 = 3 > 1$: o ponto está \emph{fora} da margem, do
lado certo. A solução do problema de margem máxima depende apenas dos pontos que
tocam a margem --- os demais têm multiplicador de Lagrange zero e poderiam ser
apagados do conjunto de dados sem alterar nada.

É a propriedade que dá nome ao método e o distingue, por exemplo, do LDA, cuja
fronteira depende das \emph{médias} e portanto de todas as observações. Mover $E$
para $(50,2)$ também não mudaria nada; mover $A$ um pouco muda tudo --- é o item
seguinte.

\smallskip
\textbf{(d)} Agora os positivos são $(2{,}5\,;1)$ e $(3,3)$, e os negativos
continuam em $x_1=1$. O ponto positivo mais próximo dos negativos é
$A=(2{,}5\,;1)$, e o negativo mais próximo dele é $C=(1,1)$: a distância entre
eles é $1{,}5$, na horizontal. O hiperplano de margem máxima é a mediatriz desse
segmento,
\[
  x_1 = \frac{1+2{,}5}{2} = 1{,}75,
  \qquad\text{margem} = \frac{1{,}5}{2} = \mathbf{0{,}75}.
\]
Na forma canônica, $\bm w = (1/0{,}75,\ 0) = (4/3,\ 0)$ e
$b = -1{,}75\cdot 4/3 = -7/3$, de modo que $1/\norm{\bm w} = 0{,}75$.

Confira que $D=(1,3)$ também toca: $(-1)(4/3 - 7/3) = 1$. Já $B=(3,3)$ dá
$(1)(4 - 7/3) = 5/3 > 1$ e ficou de fora. Os vetores de suporte são portanto
$\mathbf{A}$, $\mathbf{C}$ e $\mathbf{D}$.

Compare com o item (a): mover \emph{um} ponto de $0{,}5$ encolheu a margem em
$25\%$ e mudou o conjunto de vetores de suporte. A SVM de margem rígida é
extremamente sensível aos pontos da fronteira --- e é exatamente esse defeito que
o parâmetro $C$ do próximo exercício vem corrigir.
\end{solucao}

%------------------------------------------------------------------------------
\begin{exercicio}[o que $C$ faz, e por que ele é invertido]
A SVM de margem suave resolve
\[
  \min_{\bm w,\,b,\,\bm\xi}\ \tfrac12\norm{\bm w}^2 + C\sum_{i=1}^n \xi_i
  \qquad\text{sujeito a}\qquad
  y_i(\bm w^\top\x_i+b)\ge 1-\xi_i,\quad \xi_i\ge 0 .
\]
\begin{enumerate}[label=(\alph*),leftmargin=1.1cm]
  \item Descreva o que acontece quando $C\to\infty$. Existe sempre solução?
  \item Descreva o que acontece quando $C\to 0$. O que ocorre com $\bm w$, com a
        margem e com o número de vetores de suporte?
  \item Portanto: $C$ grande corresponde a um modelo mais \emph{flexível} ou mais
        \emph{rígido}? Compare com o papel de $\lambda$ na Ridge (Aula~02) e
        explique por que essa comparação confunde tanta gente.
  \item Ao montar uma grade de busca para $C$, por que se usa
        \texttt{np.logspace(-3, 3, 7)} e não \texttt{np.linspace(0.001, 1000, 7)}?
\end{enumerate}
\end{exercicio}

\begin{solucao}
\textbf{(a)} Com $C\to\infty$, qualquer $\xi_i>0$ paga um preço infinito, então a
solução força $\xi_i=0$ para todo $i$: é a \textbf{margem rígida} do
Exercício~1. A margem fica a mais larga possível \emph{sem que nenhum ponto a
invada}, e poucos pontos são vetores de suporte.

\textbf{Não existe sempre solução:} se os dados não forem linearmente separáveis,
o problema com $\xi_i=0$ é \emph{inviável} --- não há hiperplano que satisfaça
todas as restrições. É por isso que a margem suave existe: sem ela, o método
simplesmente não se aplica a dados reais.

\smallskip
\textbf{(b)} Com $C\to 0$ a soma dos $\xi_i$ deixa de custar, e o objetivo vira
essencialmente $\min\tfrac12\norm{\bm w}^2$, cuja solução é $\bm w\to\bm 0$.
Então:
\begin{itemize}[leftmargin=1.4cm]
  \item $\norm{\bm w}\to 0$, logo a margem $1/\norm{\bm w}\to\infty$: a faixa
        engole todo o espaço;
  \item \textbf{todos} os pontos passam a estar dentro da margem, e portanto
        todos viram vetores de suporte;
  \item a fronteira deixa de depender dos dados de forma útil --- o classificador
        degenera.
\end{itemize}

\smallskip
\textbf{(c)} $C$ \textbf{grande} é modelo \textbf{mais flexível} (menos
regularizado): ele persegue os dados de treino e aceita margens estreitas para
não errar ninguém.

A confusão com $\lambda$ da Ridge é natural, porque os dois parâmetros aparecem
em posições análogas mas com papéis \emph{opostos}. Escrevendo os dois objetivos
lado a lado:
\[
  \underbrace{\norm{\bm Y-\bm X\bbeta}^2 + \lambda\norm{\bbeta}^2}_{\text{Ridge}},
  \qquad
  \underbrace{\tfrac12\norm{\bm w}^2 + C\textstyle\sum_i\xi_i}_{\text{SVM}} .
\]
Na Ridge, o hiperparâmetro multiplica a \textbf{penalização}, então
$\lambda$ grande $\Rightarrow$ mais regularização $\Rightarrow$ mais rígido. Na
SVM, ele multiplica o \textbf{erro}, então $C$ grande $\Rightarrow$ menos
regularização $\Rightarrow$ mais flexível. Dividindo o objetivo da SVM por $C$
fica claro: $\tfrac{1}{2C}\norm{\bm w}^2 + \sum_i\xi_i$, e $1/(2C)$ é que faz o
papel de $\lambda$.

Um jeito de nunca mais errar: \textbf{$C$ é o preço de errar}. Preço alto, o
modelo não erra no treino.

\smallskip
\textbf{(d)} Porque o efeito de $C$ é \textbf{multiplicativo}, não aditivo: o que
importa é a ordem de grandeza. A diferença entre $C=0{,}001$ e $C=0{,}01$ é
enorme (um fator de 10 na regularização), e a diferença entre $C=900$ e
$C=1000$ é irrelevante.

Uma grade linear de $0{,}001$ a $1000$ com 7 pontos daria
$\{0{,}001;\ 166{,}7;\ 333{,}3;\ 500;\ 666{,}7;\ 833{,}3;\ 1000\}$: seis pontos
espremidos na região em que nada muda, e \emph{nenhum} entre $0{,}001$ e $166$ ---
justamente o intervalo onde está o ótimo na maioria dos problemas. A grade
logarítmica dá $\{10^{-3},10^{-2},\dots,10^{3}\}$, um ponto por ordem de
grandeza.

A mesma regra vale para o $\alpha$ da Ridge e do Lasso, para o $\gamma$ do kernel
RBF e para a taxa de aprendizado do \emph{boosting}.
\end{solucao}

%------------------------------------------------------------------------------
\begin{exercicio}[o truque do kernel, explicitamente]
Em $\R^2$, considere o kernel polinomial de grau 2
\[
  K(\x,\bm z) = (\x^\top\bm z + 1)^2 .
\]
\begin{enumerate}[label=(\alph*),leftmargin=1.1cm]
  \item Expanda $K$ e exiba um mapa $\varphi:\R^2\to\R^m$ tal que
        $K(\x,\bm z)=\varphi(\x)^\top\varphi(\bm z)$. Qual é $m$?
  \item Conte as operações necessárias para calcular $K(\x,\bm z)$ pelos dois
        caminhos: pela fórmula fechada e via $\varphi$.
  \item Em $\R^p$, o espaço de atributos do kernel polinomial de grau $p$ tem
        $\binom{p+p}{p}$ dimensões. Calcule para $p=100$, $p=3$ e para
        $p=1000$, $p=3$.
  \item Enuncie em uma frase o que é o ``truque do kernel'', usando os números
        do item (c).
\end{enumerate}
\end{exercicio}

\begin{solucao}
\textbf{(a)} Com $\x=(x_1,x_2)$ e $\bm z=(z_1,z_2)$,
\[
  (x_1z_1+x_2z_2+1)^2
  = x_1^2z_1^2 + x_2^2z_2^2 + 1 + 2x_1x_2z_1z_2 + 2x_1z_1 + 2x_2z_2 .
\]
Agrupando cada termo como um produto de uma parte em $\x$ por uma parte em
$\bm z$, lê-se diretamente
\[
  \varphi(\x) = \Big(x_1^2,\ x_2^2,\ \sqrt2\,x_1x_2,\ \sqrt2\,x_1,\ \sqrt2\,x_2,\ 1\Big),
\]
com $m=\mathbf{6}$. (Os $\sqrt2$ aparecem para que, ao multiplicar
$\varphi(\x)$ por $\varphi(\bm z)$, cada termo cruzado seja contado com o fator
$2$ da expansão.) Confere com $\binom{2+2}{2}=6$.

\smallskip
\textbf{(b)} Pela \textbf{fórmula fechada}: 2 multiplicações e 1 soma para o
produto interno $\x^\top\bm z$, mais 1 soma e 1 elevação ao quadrado. Cerca de
\textbf{5 operações}, e nenhuma memória extra.

Via $\varphi$: construir os dois vetores de 6 componentes (com quadrados,
produtos e multiplicações por $\sqrt2$) e depois fazer um produto interno em
$\R^6$, ou seja 6 multiplicações e 5 somas. Cerca de \textbf{25 operações}, mais
o armazenamento dos dois vetores de dimensão 6.

Em $p=2$ e $p=2$ a diferença é pequena; o item (c) mostra o que acontece quando
não é.

\smallskip
\textbf{(c)}
\[
  \binom{103}{3} = \frac{103\cdot102\cdot101}{6} = \mathbf{176\,851},
  \qquad
  \binom{1003}{3} = \frac{1003\cdot1002\cdot1001}{6} = \mathbf{167\,668\,501}.
\]
Com $p=1000$ e grau 3, o espaço de atributos tem \textbf{167 milhões} de
dimensões. Guardar um único $\varphi(\x)$ em precisão dupla custaria cerca de
$1{,}3$ GB; guardar os $\varphi(\x_i)$ de mil observações seria impossível.

\smallskip
\textbf{(d)} O truque do kernel é: \textbf{a SVM só precisa de produtos internos
entre observações, nunca das coordenadas em si} --- então dá para trabalhar num
espaço de 167 milhões de dimensões calculando, para cada par, um produto interno
em $\R^{1000}$ seguido de uma soma e um cubo, \emph{sem nunca construir esse
espaço}.

O que torna isso possível é a formulação \textbf{dual} do problema: nela $\x$
aparece exclusivamente dentro de produtos $\x_i^\top\x_j$, que podem ser
substituídos por $K(\x_i,\x_j)$. O preço é que o custo passa a ser governado por
$n$, não por $p$ --- a matriz de Gram é $n\times n$, o que explica o custo entre
$O(n^2)$ e $O(n^3)$ do treino e por que a SVM não escala para milhões de
observações.
\end{solucao}

%------------------------------------------------------------------------------
\begin{exercicio}[a perda \emph{hinge} e a esparsidade]
A SVM de margem suave é equivalente a minimizar
\[
  \sum_{i=1}^{n} \underbrace{\max\big(0,\ 1 - y_i f(\x_i)\big)}_{\text{perda \emph{hinge}}}
   \;+\; \lambda\norm{\bm w}^2,
  \qquad f(\x)=\bm w^\top\x+b .
\]
\begin{enumerate}[label=(\alph*),leftmargin=1.1cm]
  \item Complete a tabela abaixo com o valor das três perdas, em função da
        \emph{margem funcional} $m = y f(\x)$.
        \begin{center}\small
        \renewcommand{\arraystretch}{1.3}
        \begin{tabular}{@{}lccccc@{}}
        \toprule
        $m=yf(\x)$ & $-1$ & $0$ & $0{,}5$ & $1$ & $2$ \\
        \midrule
        $0$--$1$: $\1{m<0}$          & & & & & \\
        \emph{hinge}: $\max(0,1-m)$  & & & & & \\
        logística: $\ln(1+e^{-m})$   & & & & & \\
        \bottomrule
        \end{tabular}
        \end{center}
  \item Para quais valores de $m$ a perda \emph{hinge} vale exatamente zero?
        Relacione com o conceito de vetor de suporte.
  \item A perda logística vale zero para algum $m$ finito? Que consequência isso
        tem sobre \emph{quais} observações influenciam o ajuste de uma regressão
        logística?
  \item Por que não se minimiza diretamente a perda $0$--$1$, que é a que
        realmente interessa?
\end{enumerate}
\end{exercicio}

\begin{solucao}
\textbf{(a)}
\begin{center}\small
\renewcommand{\arraystretch}{1.3}
\begin{tabular}{@{}lccccc@{}}
\toprule
$m=yf(\x)$ & $-1$ & $0$ & $0{,}5$ & $1$ & $2$ \\
\midrule
$0$--$1$    & $1$ & $1$ & $0$ & $0$ & $0$ \\
\emph{hinge}& $2$ & $1$ & $0{,}5$ & $\mathbf{0}$ & $\mathbf{0}$ \\
logística   & $1{,}3133$ & $0{,}6931$ & $0{,}4741$ & $0{,}3133$ & $0{,}1269$ \\
\bottomrule
\end{tabular}
\end{center}
(Na perda $0$--$1$ o valor em $m=0$ é convenção; o ponto está exatamente sobre a
fronteira.)

\smallskip
\textbf{(b)} A \emph{hinge} vale zero exatamente quando $m\ge 1$, isto é
$y_i f(\x_i)\ge 1$: o ponto está classificado corretamente \textbf{e fora da
margem}.

Esses são precisamente os pontos que \textbf{não} são vetores de suporte. Como
eles contribuem com zero para a soma, movê-los (sem cruzar a margem) não altera
o objetivo, e portanto não altera a solução --- que é a propriedade do
Exercício~1(c), agora vista pelo lado da função de perda em vez da geometria.

A esparsidade da SVM é, portanto, consequência direta de a \emph{hinge} ter uma
região \textbf{plana em zero}. É o análogo, em perda, do que a quina do losango
$\ell_1$ faz com os coeficientes no Lasso (Aula~02): um ``bico'' na função gera
soluções exatamente nulas.

\smallskip
\textbf{(c)} \textbf{Nunca.} $\ln(1+e^{-m})>0$ para todo $m$ finito, por menor
que seja --- em $m=2$ ainda vale $0{,}1269$, e só tende a zero quando
$m\to\infty$.

A consequência é que, na regressão logística, \textbf{todas as observações
influenciam o ajuste}, inclusive as que estão longe e do lado certo da fronteira.
Não existe ``vetor de suporte'' na logística: apagar um ponto bem classificado
muda a solução, ainda que pouco. Isso a torna mais sensível a observações
distantes --- em compensação, é o que permite que ela estime probabilidades, já
que ela precisa dizer algo sobre todo o espaço, não só perto da fronteira.

\smallskip
\textbf{(d)} Porque a perda $0$--$1$ é \textbf{não convexa e descontínua}: ela é
constante por partes, com um salto em $m=0$, e sua derivada é zero em quase todo
ponto. Nenhum método baseado em gradiente tem por onde começar, e minimizá-la
exatamente é um problema NP-difícil em geral.

Tanto a \emph{hinge} quanto a logística são \textbf{cotas superiores convexas} da
perda $0$--$1$ --- confira na tabela: em cada coluna, as duas últimas linhas são
$\ge$ a primeira. Minimizar uma cota superior convexa é tratável e controla o que
interessa. É o mesmo movimento da Aula~02, em que a norma $\ell_1$ substitui a
$\ell_0$ por ser sua relaxação convexa: \emph{trocar o objetivo verdadeiro por
uma aproximação convexa que se pode de fato otimizar}.
\end{solucao}


\end{document}
